\section{Artin基底の相対window}
局所Artin環$R$上、$B$は五変数多項式環の次数別有限自由二次変形、
$W\in Z(B)_3$は滑らかな三次式へ還元するとする。
相対対角が長さ五の有限自由Koszul分解を持ち、左右で
$\operatorname{RHom}_B(R,B)=R(5)[-5]$と仮定する。
混合族では相対PBW性と有限自由Frobenius二次双対がこの仮定を与える。
$W$は左右正則で、$A=B/(W)$も次数別有限自由とする。
冪零filtrationと特殊ファイバーのNoether性から$A,B$は左右Noetherである。
\begin{theorem}\label{mix:window}
\[
 D_{\rm rel}(\operatorname{gr}A)
 =\{M\in D^b(\operatorname{gr}A):M|_B\in\Perf(\operatorname{gr}B)\},\qquad
 D_{\rm sg,rel}(A)=D_{\rm rel}/\Perf(\operatorname{gr}A)
\]
と置く。有限自由graded MFの冪等完備化は$D_{\rm sg,rel}(A)$と$R$線形Morita同値である。
実際の斉次Ore局所化による構造代数層を$\mathcal A$とすると、同じwindowと二つのmutationにより
\[
 D_{\rm sg,rel}(A)\simeq
 \mathcal K(\mathcal A)=\langle\mathcal O_{\mathcal A},\mathcal O_{\mathcal A}(1)\rangle^\perp
 \subset\Perf(\mathcal A)
\]
となる。比較は特殊ファイバーの標識およびArtin基底間の導来基底変換と両立する。
基底自身の特異性を含む絶対特異圏との同一視ではない。
\end{theorem}
\begin{proof}
相対対象は、$R$上一様な有限Tor振幅を持ち、各内部次数が$R$上perfectな対象に等しい。
一方向は有限$B$自由分解から従う。
逆に有限$A$自由分解を十分下で切ると、有限Tor振幅から高いsyzygyは$R$平坦となる。
$B$の有限相対対角分解をtensorして長さ五の$B$射影分解を得る。
ここで全内部次数の直和が有限階数であることは要求しない。

MFの$P$は特殊ファイバーで単射であり、各内部次数の単元minorから余核は$R$自由、$B$射影次元高々一となる。
逆に$R$平坦な有限生成$A$加群$M$は$\operatorname{pd}_BM\le5$であり、
$A$自由加群の$B$射影次元が一なので
$\operatorname{pd}_B\Omega_AM\le\max\{1,\operatorname{pd}_BM-1\}$である。
十分高いsyzygy$M'$の分解$0\to P_1\xrightarrow P P_0\to M'\to0$を取る。
$WM'=0$から$W:P_0(-3)\to P_0$は一意に$P$を通り$Q$を与える。
$PQ=W$と$P$の単射性から$QP=W$も従う。
有限graded射影は局所基底上の冪等元持上げにより自由なshiftの直和因子として扱える。
一般の相対複体も有限$A$自由複体を除けばこの状況になる。
二周期完全分解とその双対の完全性は$B$へ持ち上げて$W$の正則性を使えば得られ、
安定射、MFのhomotopy類、相対特異圏の射を同一視する。
これが完全忠実性と本質的全射性を与える。

windowに必要な内部cutoffを確かめる。
高い自由分解は上の二周期tailに一致し、二段ごとに生成次数が三増すので、
任意の$i$より小さい生成次数は十分左には現れない。
内部切断$M_{\ge i}$とその商は各次数の直和因子で有限Tor振幅を保つ。
有限内部次数の相対対象は$R(j)\otimes_R\Perf(R)$で生成される。
通常のcohomological truncationで相対部分圏が閉じることは仮定しない。
中心元の長さ一分解から
\[
 \operatorname{RHom}_B(M,B)=\operatorname{RHom}_A(M,A)(3)[-1],\qquad
 \operatorname{RHom}_A(R,A)=R(2)[-4]
\]
となる。従って$D_A=\operatorname{RHom}_A(-,A)$は相対部分圏を保ち$D_A^2=1$、
$D_A(R(j))=R(2-j)[-4]$, $D_A(A(j))=A(-j)$である。
これは相対Gorenstein parameter二の式であり、Artin環の絶対有限大域次元を用いない。

Ore chartの代数$C$と$C^e$は$R$平坦である。
特殊ファイバーの滑らかなaffine対角の有限射影分解を有限自由部分分解として持ち上げる。
各核は$R$平坦で、最後の核の剰余体還元は射影である。
冪零イデアルを越えた冪等元持上げと
$\operatorname{Tor}^{C^e}_1(C^e/\mathfrak{m}_RC^e,K)
=\operatorname{Tor}^R_1(R/\mathfrak{m}_R,K)=0$
から最後の核も射影となる。
よってchartは相対smoothで、$R$上有限Tor振幅のcoherent複体は正確にchart上perfectである。
有限affine coverのČech複体は平坦な項と一様有限振幅を持つ。
proper特殊ファイバーの有限cohomologyと冪零摂動の有限和を用いて、各twistの導来切断は$R$上perfectで基底変換と両立する。
Serre消滅と冪零filtrationにより$\mathbf R\Gamma_{\ge i}$は有限生成かつ一様有限振幅で、相対部分圏に入る。

Burke--Stevenson\cite{BS}の原定理は次数零部分の有限大域次元を仮定するため、Artin基底へ直接適用しない。
以上の内部cutoff、torsion局所化、perfect双対性、$R(j),A(j)$の像を用いて内部次数分解の直交補を比較する。
parameter二なので二つの$A$-shiftが残り、window $i=-1$は
\[
 \Perf(\mathcal A)=\langle\mathcal O_{\mathcal A},\mathcal O_{\mathcal A}(1),\mathcal B_{-1}\rangle,
 \qquad\mathcal B_{-1}\simeq D_{\rm sg,rel}(A)
\]
を与える。この順序の最後の成分をそのまま$\mathcal K$と呼ばず、
$L_{\mathcal O}L_{\mathcal O(1)}$で左へ移して
$\Perf(\mathcal A)=\langle\mathcal K,\mathcal O,\mathcal O(1)\rangle$とする。
評価複体はperfectなのでmutationも有限である。
全構成は自由分解、導来随伴、cone、dg quotient上で行われ、Morita同値へ上がる。
有限$B$分解、Ore局所化、perfect評価複体を用いるので導来基底変換と両立する。
両側で同じwindowとmutationを用いれば標識も一致する。
\end{proof}

\begin{lemma}[生成子一つの持上げによる基底変換]
$Q$は有限大域次元の可換Noether局所環、$R$はその局所Artin商とする。
次数別有限自由な$(B_Q,W_Q)$について、有限自由MFの有限直和$G_Q$があり、
$G_k$が剰余体MF圏を生成し、$A_Q=\operatorname{RHom}(G_Q,G_Q)$が$Q$上smooth properであるとする。
剰余体MF圏もproperとする。このとき
\[
 \Perf(A_Q\otimes_Q^{\mathbf L}R)\simeq\MF(B_R,W_R)^{\rm perf}.
\]
任意のArtin対象そのものの$Q$への持上げは必要ない。
\end{lemma}
\begin{proof}
有限自由graded MFの射複体は各cohomological次数が有限自由で、剰余体還元は実際の射複体である。
その剰余体複体を有限cohomologyとcontractibleな二項複体へ分け、基底を$R$へ持ち上げる。
微分の冪零摂動は有限幾何級数で消去できるので、射複体は$R$上perfectとなる。
各次数は有限自由で、摂動合成は基底の冪零次数で止まるため、一方向無限の複体でもこの操作は定義される。
$Q$上の$R$の有限射影分解により$A_Q\otimes_Q^{\mathbf L}R=A_R$を得る。
任意の$M$に対して$P=\operatorname{RHom}(G_R,M)$は$R$上perfectであり、
$A_R$のsmooth対角の有限cell分解をtensorすると$P$は$A_R$上perfectになる。
従って評価$P\otimes_{A_R}^{\mathbf L}G_R\to M$のcone $Z$は$G_R$と直交する。
導来基底変換と剰余体生成性から$Z_k=0$、perfect複体のNakayamaから$\operatorname{RHom}(Z,Z)=0$、従って$Z=0$である。
\end{proof}
